\section{Outlier-Robust Convex Segmentation}\label{sec:outlierRobustObj}
Segmentation is the task of dividing a sequence of $n$ data samples $\{\vxii\}_{i=1}^n$, into $K$ groups of consecutive samples, or segments, such that each group is homogeneous with respect to some criterion. A common choice of such a criterion often involves minimizing the squared Euclidean distance of a sample to some representative sample $\vmuii$. %, e.g. the centroid of the segment.
This criterion is highly sensitive to outliers and indeed, as we show empirically below, the performance of segmentation algorithms degrades drastically when the data is contaminated with outliers.
It is therefore desirable to incorporate robustness to outliers into the model. We achieve this by allowing some of the input
samples $\vxii$ to be identified as outliers, in which case we do not require $\vmuii$ to be close to these samples.
To this end we propose to minimize:
\begin{align*}%\label{eq:outlierRobustObjective}
    \min_{\vmu,\vz}\bbraces{&\frac{1}{2}\sum_{i=1}^n\normt{\vxii-\vzii-\vmuii}} \\
    \textrm{ s.t. } & \sum_{i=1}^{n-1}\ind{\vnorm{\vmui{i+1}-\vmuii}_p>0}\leq K-1~,\\ %\sum_{i=1}^{n-1}\vnorm{\vmui{i+1}-\vmuii}_2\\&+
                    & \sum_{i=1}^n\ind{\vnorm{\vzii}_q>0}\leq M~,
\end{align*}
where $p,q\geq 1$. Considering samples $\vxii$ for which $\vzii=0$, the objective measures the loss of replacing a point $\vxii$ with some shared point $\vmuii$, and can be thought of as minus the log-likelihood under Gaussian noise. Samples $i$ with $\vzii\neq0$ are intuitively identified as outliers. The first constraint bounds the number of segments by $K$,
%and is the $\ell_0$ norm of the vector of norms-of-differences of the $\vmuii$s.
while the second constraint bounds the number of outliers by $M$.
%and can be seen as the $\ell_0$ norm of the vector with $\vnorm{\vzii}_q$ as its $i$th component.
The optimal value for a nonzero $\vzii$ is to set $\vzii=\vxii-\vmuii$,
making the contribution to the objective zero, and thus in practice
ignoring this sample, treating it as an outlier. We note that a similar approach to robustness was employed by \cite{forero2011outlier} in the context of robust clustering, and by \cite{mateos2012robust} in the context of robust PCA.
Since the $\ell_0$ constraints results in a non convex
problem, we use a common practice and replace it with a convex surrogate $\ell_1$ norm which induces
sparsity. For the $\vmuii$ variables it means that for most samples we will have $\vmui{i+1}-\vmui{i}=0$, allowing the identification of the corresponding samples as belonging to the same segment.
For the $\vzii$ variables, it means that most will satisfy $\vzii=0$ and for some of them, the
outliers, otherwise. We now incorporate the relaxed constraints into the objective, and in addition consider a slightly more general formulation in which we allow weighting of the summands in the first constraint. We get the following optimization problem:
\begin{align}\label{eq:outlierRobustObjective}
    \min_{\vmu,\vz}\bbraces{ &\frac{1}{2}\sum_{i=1}^n\normt{\vxii-\vzii-\vmuii}+\lambda\sum_{i=1}^{n-1}w_i\vnorm{\vmui{i+1}-\vmuii}_p\\
    &+\gamma\sum_{i=1}^n\vnorm{\vzii}_q}~,\nonumber
\end{align}
where $w_i$ are weights to be determined.
The parameter $\lambda>0$ can be thought of as a tradeoff parameter between the first term which is minimized with $n$ segments, and the second term which is minimized with a single segment. As $\lambda$ is decreased, it crosses values at which there is a transition from $K$ segments to $K+1$ segments, in a phase-transition like manner where $1/\lambda$ is the analog of temperature. The parameter $\gamma>0$ controls the amount of
outliers, where for $\gamma=\infty$ we enforce $\vzii=0$ for all samples,
and for $\gamma=0$ the objective is optimal for $\vzii=\vxii-\vmuii$, and thus all
samples are in-fact outliers.
Alternatively, one can think of $\lambda,\gamma$ as the Lagrange multipliers of a constrained optimization problem. In what follows we consider $p=2$ and $q=1,2$, focusing empirically on $q=2$. Note that $q=1$ encourages sparsity of coordinates of $\vzii$, and not of the vector as a whole.
%, as is the case for $q=2$.
This amounts to outliers being modeled as noise %concentrated
in few features or samples, respectively.

%Specifically, we introduce
%additional variables $\{\vzii\}$ to the objective in
%\eqref{eq:segmentationObjective}, which, intuitively, are
%identified as the perturbation or outlier values, and replace
%\eqref{eq:segmentationObjective} with the following,
%\begin{align*}%\label{eq:outlierRobustObjective}
%    \min_{\vmu,\vz}\bbraces{&\frac{1}{2}\sum_{i=1}^n\normt{\vxii-\vzii-\vmuii}+\lambda\sum_{i=1}^{n-1}\vnorm{\vmui{i+1}-\vmuii}_2\\&+
%    \gamma\sum_{i=1}^n\ind{\vnorm{\vzii}_q>0}}~.%\nonumber
%\end{align*}
%When $\vzii=\vzero$ for all $i=1,\ldots,n$, \eqref{eq:outlierRobustObjective} reduces to
%\eqref{eq:segmentationObjective}. Ideally, as before, the last term counts and therefore regularizes the number of samples for which $\vzii\neq 0$, and can be seen as the $\ell_0$ norm of the vector with $\vnorm{\vzii}_q$ as its $i$th component. In this case the optimal value for a nonzero $\vzii$ is to set $\vzii=\vxii-\vmuii$,
%making the contribution to the first term zero, and thus in practice
%ignoring this sample, treating it as an outlier.
%However, as before, the $\ell_0$ regularization results in a non convex
%function, and thus we relax it by an $\ell_1$ norm, which encourages
%that for most samples we have $\vzii = 0$ and for some of them, the
%outliers, otherwise. The parameter $\gamma$ controls the amount of
%outliers. For $\gamma=\infty$ we enforce $\vzii=0$ for all samples,
%reducing to the algorithm of previous section, while for $\gamma=0$
%the objective is optimal for $\vzii=\vxii-\vmuii$, and thus all
%samples are in-fact outliers.
%In what follows we consider $q=1,2$, noting that $q=1$ encourages sparsity of individual components of $\vzii$, and not on the vector as a whole, as is the case for $q=2$. The resulting convex optimization problem is given by
%\begin{align}\label{eq:outlierRobustObjective}
%    \min_{\vmu,\vz}\bbraces{&\frac{1}{2}\sum_{i=1}^n\normt{\vxii-\vzii-\vmuii}+\lambda\sum_{i=1}^{n-1}\vnorm{\vmui{i+1}-\vmuii}_2\\&+
%    \gamma\sum_{i=1}^n\vnorm{\vzii}_q}~.\nonumber
%\end{align}

\subsection{Algorithms}\label{sec:algorithms}
The decoupling between $\vmu$ and $\vz$ allows us to optimize
\eqref{eq:outlierRobustObjective} in an alternating manner, and we call this algorithm %for alternately optimizing \eqref{eq:outlierRobustObjective}
Outlier-Robust Convex Sequential (ORCS) segmentation.
Holding $\vmu$ constant, optimizing
over $\vz$ is done analytically by noting that
\eqref{eq:outlierRobustObjective} becomes the definition of the
proximal operator evaluated at $\vxii-\vmuii$, for which a closed-form solution exists.
For $q=1$ the objective as a function of $\vz$
is separable both over coordinates and over data samples, and the proximal
operator is the shrinkage-and-threshold operator evaluated at each
coordinate $k$:
\begin{equation*}%\label{eq:proxOp_l1}
    \mathrm{prox}_\gamma\paren{v_k}=\sign\paren{v_k}\cdot\max\braces{0,\abs{v_k}-\gamma}.
    %\begin{cases}
%        v_k-\gamma, & v_k\geq\gamma\\
%        v_k+\gamma, & v_k\leq-\gamma\\
%        0, & \mathrm{otherwise}~.
%    \end{cases}
\end{equation*}
However, we are interested in zeroing some of the $\vzii$'s as a whole, so we set $q=2$. In this case, the objective is separable over data samples, and the proximal operator is calculated to be:
\begin{align}\label{eq:proxOp_l2}
%\mathrm{prox}_{\gamma,\ell_2}\paren{v} &=&
\mathrm{prox}_{\gamma}\paren{\vv} = \vv\cdot\max\braces{0,1-\frac{\gamma}{\vnorm{\vv}_2}}.
 % \begin{cases}
%    %\paren{1+\frac{\gamma}{\vnorm{v}_2-\gamma}}^{-1}v, &\text{if $\vnorm{v}_2>\gamma$}\\
%    \frac{\vnorm{v}_2-\gamma}{\vnorm{v}_2}v, &\text{if $\vnorm{v}_2>\gamma$}\\
%    0, & \text{otherwise}~.
%  \end{cases}
\end{align}
Holding $\vz$ constant, optimizing over $\vmu$ is done by defining
$\hvxii\triangleq \vxii-\vzii$, which results in the following optimization problem:
\begin{equation}\label{eq:segmentationObjective_mu}
 \min_\vmu\sum_{i=1}^n\normt{\hvxii-\vmuii}  + \lambda\sum_{i=1}^{n-1}w_i\vnorm{\vmui{i+1}-\vmuii}_2.
\end{equation}
%We note that in carrying the alternating optimization, using the optimal $\vmu$ from one iteration as the initial $\vmu$ for the next iteration speeds-up the convergence significantly.
Note that \eqref{eq:segmentationObjective_mu} is equivalent to \eqref{eq:outlierRobustObjective} with no outliers present. We also note that if we plug the analytical solution for the $\vzii$s into \eqref{eq:segmentationObjective_mu} (via the $\hvxii$s), the loss term turns out to be the multidimensional equivalent of the Huber loss of robust regression.
We now discuss two approaches for solving \eqref{eq:segmentationObjective_mu}, either exactly or approximately.

\paragraph{Exact solution of \eqref{eq:segmentationObjective_mu}:}
The common proximal-gradient approach \cite{bach2011convex,Beck_FISTA:2009} for solving non-smooth convex problems has in this case the disadvantage of convergence time which grows linearly with the number of samples $n$. The reason is that the Lipschitz constant of the gradient of the first term in \eqref{eq:segmentationObjective_mu} grows linearly with $n$, which results in a decreasing step size. An alternative approach is to derive the dual optimization problem to \eqref{eq:segmentationObjective_mu}, analogously to the derivation of \cite{Beck_TV:2009} %%(see also \cite{Chambolle:2004})
in the context of image denoising. The resulting objective is smooth and has a bounded Lipschitz constant independent of $n$.
Yet another approach was proposed by \cite{bleakley2011group} for the task of change-point detection, who showed that under a suitable change of variables \eqref{eq:segmentationObjective_mu} can be formulated as a group-LASSO regression \cite{tibshirani1996regression,yuan2006model}.
%see also \cite{levy2007catching} for the one-dimensional case.
%They proposed two algorithms for optimizing \eqref{eq:segmentationObjective_mu} either exactly or approximately.
\paragraph{Approximate solution of \eqref{eq:segmentationObjective_mu}:}\label{par:seg_approx_solution}
Two reasons suggest that deriving an alternative algorithm for solving \eqref{eq:segmentationObjective_mu} might have an advantage. %over the exact solution.
First, the parameter $\lambda$
does not allow direct control of the resulting number of segments, and in many use-cases such a control is a desired property.
%The second reason has theoretical roots.
Second, as mentioned above,
%\cite{vert2010fast}
\cite{bleakley2011group} showed that \eqref{eq:segmentationObjective_mu} is equivalent to group-LASSO regression, under a suitable change of variables. It is known from the theory of LASSO regression
that certain conditions on the design matrix %(the matrix of inputs, or predictors, for the regression)
must hold in order for perfect detection of segment boundaries to be possible. Unfortunately, these conditions are violated for the objective in \eqref{eq:segmentationObjective_mu} %when it is formulated as group-LASSO regression
; see \cite{levy2007catching} and references therein. Therefore a non-exact solution has a potential of performing better, at least in some situations. We indeed encountered this phenomenon empirically, as is demonstrated in \secref{sec:experimentalResults}.
Therefore we also derive an alternative top-down, greedy algorithm, which finds a segmentation into $K$
segments, where $K$ is a user-controlled parameter. The algorithm works in rounds. On each round it picks a segment of a
current segmentation, and finds the optimal segmentation of it into two subsequences.
We start with the following lemma, which gives an analytical rule which solves \eqref{eq:segmentationObjective_mu} for the case of $K=2$ segments.
\begin{lemma}\label{lem:analyticalSolutionK=2}
Consider the optimal solution of \eqref{eq:segmentationObjective_mu} for the largest parameter $\lambda$ for which there are $K=2$ segments, and denote this value of the parameter by $\lambda^*$. Denote by $i^*$ the associated splitting point into $2$ segments, i.e. samples $\hvxii$ with $i\leq i^*$ belong to the first segment, and otherwise belong to the second segment. Then %$i^*$ is given by
$i^*\paren{\vx} =\argmax{1\leq i\leq n-1}g\paren{i,\vx}$, where:
\begin{align}\label{eq:lambdaStarWeighted}
    %i^*\paren{\vx} &=\argmax{1\leq i\leq n-1}g\paren{i,\vx},\\% \textrm{ and }
    %\lambda^*\paren{\vx} = g\paren{n^*,\vx},\\
    g\paren{i,\vx} = \bbraces{\frac{i(n-i)}{w_in}\vnorm{\bvxi{2}(i)-\bvxi{1}(i)}_2},%\nonumber
\end{align}
and $\bvxi{1,2}(i)$ are the means of the first and second segments, respectively, given that the split occurs after the $i$th sample. In addition, $\lambda^*\paren{\vx} = g\paren{i^*,\vx}$.
\end{lemma}
The proof is given in %\secref{sec:proof_lem:analyticalSolutionK=2} of
the supplementary material.
%of \cite{DBLP:journals/corr/LattimoreCS14}.
This result motivates a top-down hierarchical segmentation algorithm, which chooses at each iteration to split the segment
which results in the maximal decrease of the sum-of-squared-errors criterion. Note that we
cannot use the criterion of minimal increment to the objective in \eqref{eq:segmentationObjective_mu}, since by continuity of
the solution path, there is no change in the objective at the splitting from $K=1$ to $K=2$ segments. The top-down algorithm can be implemented in $\mathcal{O}\paren{nK}$.
%, where $n$ is the number of data samples and $K$ the number of segments.
It has the advantage that no search in the solution path is needed in case $K$ is known, and
that this search can be made efficiently in case where $K$ is not known. The top-down approach is used in the algorithm presented in \secref{sec:robust_top_down}.

From the functional form of $g\paren{i,\vx}$ in \eqref{eq:lambdaStarWeighted} it is clear that in the unweighted case ($w_i=1$ for all $i$), the solution is biased towards segments of approximately the same length, because of the $i\paren{n-i}$ factor.
We now show that a specific choice of $w_i$ exactly recovers the solution to the unrelaxed optimization problem, where the regularization term in \eqref{eq:segmentationObjective_mu} is replaced with the $\ell_0$ constraint, that is $\sum_{i=1}^{n-1}\ind{\vnorm{\vmui{i+1}-\vmuii}_2>0}=1$. This is formulated by the following lemma:
\begin{lemma}\label{lem:weighted_solution_l0_l1}
    Consider the case of two segments, $K=2$, and denote by $n_w^*$ the minimizer of \eqref{eq:lambdaStarWeighted} with $w_i=\sqrt{i\paren{n-i}}$. Then the split into two segments found by solving the following: % optimization problem:
    \begin{align}\label{eq:segmentationObjective_unrelaxed_K=2}
        \argmin{\vmu}&\braces{\sum_{i=1}^n\normt{\vxii-\vmuii}},\\
        \mathrm{s.t.}\ &\sum_{i=1}^{n-1}\ind{\vnorm{\vmui{i+1}-\vmuii}_2>0}=1~,\nonumber
    \end{align}
    is also given by $n_w^*$.
\end{lemma}
The proof appears in %\secref{sec:proof_lem:weighted_solution_l0_l1} of
the supplementary material.
%of \cite{DBLP:journals/corr/LattimoreCS14}.
We note that the same choice for $w_i$ was derived by \cite{bleakley2011group} from different considerations based on a specific noise model for the stochastic process generating the data. In this sense our derivation is more general, as it does not make any assumptions about the data.

\subsubsection{Robust top-down algorithm}\label{sec:robust_top_down}
We now propose a robust top-down algorithm for approximately optimizing \eqref{eq:outlierRobustObjective}.
For a fixed value of $\vmu$, using \eqref{eq:proxOp_l2} we can
calculate analytically which of the $\vzii$s represent a detected
outlier. These are $\vzii$s which satisfy
$\vnorm{\vzii}_2>\gamma$. This allows us to calculate the value
$\gamma^*$ for which the first outlier is detected as having a
non-zero norm. Furthermore, for $\lambda=\lambda^*,\ \gamma=\gamma^*$ we know that $\vzii=0$ for all $i=1\comdots n$, and therefore
$\vmuii=\mathrm{mean}\paren{\vx-\vz}=\bvx$ for all $i=1\comdots n$, and we
can find $\gamma^*$ analytically:
\begin{equation}\label{eq:gmaStar}
    \gamma^*=\max_i\braces{\vnorm{\vxii-\bvx}_2}~,
\end{equation}
where the index $i^*$ at which the maximum is attained is the index to the first detected outlier. The value of $\lambda^*$ is found as given in \lemref{lem:analyticalSolutionK=2}, with the replacement of each $\vxii$ with $\hvxii$ as defined above. We note that the values $\lambda^*,\ \gamma^*$ are helpful for finding a solution path, since they allow to exclude parameters which result in trivial solutions.

In the case where $\lambda=\lambda^*$ we can extend \eqref{eq:gmaStar} for any number $M>1$ of outliers, by simply looking for the first $M$ vectors $\vxii-\bvx$ with the largest norm.
%\begin{equation}\label{eq:gmaStar}
%    \gamma^*_M=\vnorm{\vxii-\bvx}_2\textrm{: }\vxii \textrm{ is the vector with the $M$th largest $\ell_2$ norm}
%\end{equation}
In this case it no longer holds true that $\vzii=0$ for all $i=1\comdots n$, so
we have to use the alternating optimization in order to find a
solution. However, each iteration is now solved analytically and
convergence is fast compared to the case $\lambda<\lambda^*$ where we
do not have an analytical solution for the optimization over $\vmu$. This
result motivates the top-down version of the ORCS algorithm. We
denote the algorithm by TD-ORCS for the unweighted case ($w_i=1$, $i=1\comdots n$), and by WTD-ORCS when using the weights given in \lemref{lem:weighted_solution_l0_l1}. The number of required segments $K$ and number of
required outliers $M$ is set by the user. In
each iteration the algorithm chooses the segment-split which results in
the maximal decrease in the squared loss. Whenever a segment is split,
the number of outliers belonging to each sub-segment is kept and used
in the next iteration, so the overall number of outliers equals $M$ at
all iterations. The algorithm is summarized in \algref{alg:topdownORCS}.

\begin{algorithm}[!t]
\footnotesize
 \caption{Top-down outlier-robust hierarchical segmentation}\label{alg:topdownORCS}
  \begin{algorithmic}
    \STATE \textbf{Input:} Data samples $\braces{\vxii}_{i=1}^n$.
    \STATE \textbf{Parameters:} Number of required segments $K$, number of required outliers $M$, weights $\braces{w_i}_{i=1}^{n-1}$.
    \STATE \textbf{Initialize:} Segments boundaries $\mathrm{B}=\braces{1,n}$, current number of segments $k=\abs{B}-1$, number of outliers for each segment $M_j=M$ (for $j=1$).
    \WHILE{$k<K$}
        \FOR{$j=1,...,k$}
            \STATE  Set $S_j=\braces{\vxi{B_j},..,\vxi{B_{j+1}-1}}$.
            \WHILE {not converged}
                \STATE Split segment $S_j$ into sub-sequences $S_{1,2}$, using the solution to
                \eqref{eq:lambdaStarWeighted} at $\vx-\vz$, with weights $\braces{w_i}$.
                \STATE Find $\gamma$ for $M_j$ outliers, using the extension of \eqref{eq:gmaStar} to $M_j$ outliers.
                \STATE Set $\vzii=\mathrm{prox}_{\gamma}\paren{\vxii-\bvx}$, for $i=1\comdots n$.
            \ENDWHILE
            \STATE Calculate the mean $\bvxi{j}$ of segment $S_j$, and the means of $S_{1,2}$, denoted by $\bvxi{1}$ and $\bvxi{2}$.
            %\item Set $L(j)=\sum\limits_{i\in S_1}\normt{\vxii-\bvxi{1}}+\sum\limits_{i\in S_2}\normt{\vxii-\bvxi{2}}-\sum\limits_{i\in S_j}\normt{\vxii-\bvxi{j}}$.
            \STATE Set \(L(j)=%\negspace
              \sum \limits_{i\in S_1}\!\normt{\vxii-\bvxi{1}}
              \!\!+\!\!
              \sum\limits_{i\in S_2}\!\normt{\vxii-\bvxi{2}}
              \!\!-\!\!
              \sum \limits_{i\in S_j}\!\normt{\vxii-\bvxi{j}}\).
        \ENDFOR
        \STATE Choose segment $j^* = \arg\max_j  L(j)$ and calculate the splitting
        point $n_{j^*}$
        \STATE Add the new boundary $\vxi{B_{j^*}} + n_{j^*} + 1$ to the (sorted) boundary list $B$.
        \STATE Set $M_j$ and $M_{j+1}$ to the number of $\vzii$s with non-zero norm in segment $S_j$ and $S_{j+1}$, respectively.
    \ENDWHILE
\end{algorithmic}
\end{algorithm}

%\begin{algorithm}[!t]
%\footnotesize
% \caption{Top-down outlier robust hierarchical segmentation}\label{alg:topdownORCS}
%  \begin{algorithmic}
%    \STATE \textbf{Input:} Data samples $\vxii\in\mathbb{R}^d$, $i=1 \comdots n$.
%    \STATE \textbf{Parameters:} Number of required segments $K$, number of required outliers $M$.
%    \STATE \textbf{Initialize:} Segments boundaries $\mathrm{B}=\braces{1,n}$, current number of segments $k=\abs{B}-1$, number of outliers for each segment $M_j=M$ (for $j=1$).
%    \WHILE{$k<K$}
%        \FOR{$j=1,...,k$}
%        \STATE
%        \begin{itemize}
%          \nolineskips
%            \item Set $S_j=\braces{\vxi{B_j},..,\vxi{B_{j+1}-1}}$
%            \item {\bf While} {not converged} {\bf do}
%            \begin{itemize}
%              \nolineskips
%                \item Find split for segment $S_j$ to sub-sequences $S_{1,2}$, by evaluating \eqref{eq:lambdaStar}, or
%                    \eqref{eq:lambdaStar_weighted} for the weighted version, at $\vx-\vz$.
%                \item Find $\gamma$ for $M_j$ outliers, using the extension of \eqref{eq:gmaStar} to $M_j$ outliers.
%                \item Set $z=\mathrm{prox}_{\gamma}\paren{\vx-\bvx}$
%                  \end{itemize}
%
%            \item Calculate the mean $\bvxi{j}$ of segment $S_j$, and the means of $S_{1,2}$, denoted by $\bvxi{1}$ and $\bvxi{2}$.
%            %\item Set $L(j)=\sum\limits_{i\in S_1}\normt{\vxii-\bvxi{1}}+\sum\limits_{i\in S_2}\normt{\vxii-\bvxi{2}}-\sum\limits_{i\in S_j}\normt{\vxii-\bvxi{j}}$.
%            \item Set \(L(j)\!\!=\!\!
%              \sum \limits_{i\in S_1}\!\normt{\vxii-\bvxi{1}}
%              \!\!+\!\!
%              \sum\limits_{i\in S_2}\!\normt{\vxii-\bvxi{2}}
%              \!\!-\!\!
%              \sum \limits_{i\in S_j}\!\normt{\vxii-\bvxi{j}}\).
%        \end{itemize}
%        \ENDFOR
%        %        \State Choose segment $j^*=\argmax{j}L(j)$ and calculate its splitting point $n_{j^*}$.
%        %      \State Add the new boundary $\vxi{B_{j^*}}+n_{j^*}+1$ to the (sorted) boundary list $B$.
%        \STATE Choose segment $j^* \!\! = \!\! \argmax{j}  L(j)$ and calculate the splitting point $n_{j^*}$
%        \STATE Add the new boundary $\vxi{B_{j^*}} \! + \! n_{j^*} \! + \! 1$ to the (sorted) boundary list $B$
%        \STATE Set $M_j$ and $M_{j+1}$ to the number of $\vzii$'s with non-zero norm in segment $S_j$ and $S_{j+1}$, respectively
%    \ENDWHILE
%\end{algorithmic}
%\end{algorithm}
%
